\section{Analysis}
\label{sec:analysis}

This section describes what we saw when training UltraBERT.
Unless noted, the numbers are from the dev set with one seed per setting.

\subsection{Learning rate and training length}

The learning rate that is best on dev, $1.5\cdot10^{-5}$ for both swept recipes, is also among the best on the external benchmarks.
Higher rates add little on dev and lower the external accuracy, down to 53.2\% for the softmax at $5\cdot10^{-5}$.
Qwen3-0.6B shows the same pattern more strongly: its external accuracy falls with every increase of the rate, from 57.4\% at $5\cdot10^{-6}$ to 53.9\% at $2\cdot10^{-5}$.
A second epoch has the same effect.
With the softmax at $3\cdot10^{-5}$ it raises dev accuracy from 70.7\% to 73.9\%, above any one-epoch encoder run, and lowers external accuracy from 56.2\% to 55.1\%.
We report one epoch for all models, since a longer run fits the mixture at the cost of behaviour outside it.

\subsection{Datasets}

\Cref{tab:perdataset} in the appendix gives the test accuracy for every dataset, averaged over seeds.
The sigmoid-normalised encoder is above the fully fine-tuned decoder on 10 of the 13 datasets.
The gaps are largest on CosmosQA, ANLI and QuALITY, where it is ahead by 15.3, 9.9 and 9.8 points, and on MASSIVE (8.0) and MuSiQue (6.9).
It is below the decoder on three datasets: Nimble ($-6.6$), LogiQA~2.0 ($-2.2$) and ProofWriter ($-1.7$).
The 4B decoder is ahead of the encoder on all ten datasets outside intent routing.
The gap is 12 to 19 points on ReClor, Nimble and LogiQA~2.0, and below one point on CosmosQA, ANLI and ContractNLI.
On the three intent datasets, CLINC150, MASSIVE and Bitext, the encoder is level with or slightly above the 4B decoder, by 1.4 to 2.1 points.

\subsection{Loss and score}

\Cref{tab:ablation} separates the changes between the plain softmax and the sigmoid-normalised loss.
Scoring with the difference of the ``yes'' and ``no'' logits, with the softmax kept, raises the external accuracy by 2.0 points and the dev accuracy by 0.4 points; averaged over seeds, the score accounts for the whole external gain (\CpExtYes{}\%, \CpExtYesNo{}\%, \CpExtSig{}\% in \cref{tab:coupling}).
Replacing the softmax with the sigmoid-normalised loss then raises dev accuracy by a further 1.3 points, and MuSiQue most of all, from 64\% to 76\%.
Nimble falls from 55\% with the plain score to 43\% with the sigmoid-normalised loss in this one run, on 100 dev rows.
On the test set, averaged over seeds, the sigmoid-normalised model is instead 3.7 points above the softmax on Nimble.
Replacing ``yes'' with ``correct'' did not help in this run: 70.6\% against 71.1\% with the ``yes'' logit, and 70.1\% against 71.5\% for the difference with ``incorrect''.

The recipes also predict different labels.
On MuSiQue the correct answer is ``Yes'' in 48\% of the rows.
The sigmoid-normalised models predict ``Yes'' for 47 to 60\% of the rows, against 64\% and 69\% for the two softmax variants (dev rows).
On ProofWriter, where ``Unknown'' is correct for a third of the rows, the softmax never predicts it and the sigmoid-normalised models predict it for 2 to 12\%.
This class stays largely unlearned for every variant.

\begin{table}[t]
  \centering\small
  \caption{Encoder variants on the dev rows, one seed (accuracy, \%). The first row is the softmax with the ``yes'' logit,
  the setting used for the softmax results elsewhere. Bold: best in each column.}
  \label{tab:ablation}
  % Generated by paper/figures/make_figures.py; do not edit.
\begin{tabular}{@{}lccccc@{}}
\toprule
Variant & Dev ALL & External & MuSiQue & ContractNLI & Nimble \\
\midrule
Softmax, ``yes'' & 71.1 & 57.1 & 57.0 & 75.0 & 55.0 \\
Softmax, ``yes'' $-$ ``no'' & 71.5 & 59.1 & 64.0 & 78.0 & 48.0 \\
Sigmoid-norm, ``yes'' $-$ ``no'' & 72.8 & 59.4 & 76.0 & 81.0 & 43.0 \\
Softmax, ``correct'' & 70.6 & 56.8 & 48.0 & 77.0 & 52.0 \\
Softmax, ``correct'' $-$ ``incorrect'' & 70.1 & 58.0 & 54.0 & 76.0 & 46.0 \\
\bottomrule
\end{tabular}

\end{table}

\subsection{Numerical precision}
\label{sec:precision}

The encoder has to be scored with fp32 weights under bf16 autocast.
Its masked-token logits are of the order of 200, where bf16 can only step by 1 or 2, so options tie or swap.
With bf16 weights the MuSiQue accuracy drifts by up to 8 points.
The decoders are scored with bf16 weights, and re-scoring a Qwen checkpoint in another precision moved its dev accuracy by at most 0.8 points.
